\section{The Bailout Protocol}\label{sec:bailout}

The Bailout Protocol is the mandatory exception-escalation mechanism of the \bxthree{} Framework. It is not an error-handling routine selected at runtime, nor an optional escalation pathway available to a system that chooses to invoke it. It is an architectural compulsion: any unresolved exception state encountered by any node in a \bxthree{} system \emph{must} propagate upward through the accountability chain until it reaches a human principal, and no machine actor along that chain may absorb, resolve, or suppress the exception. This section provides the complete formal specification: the deterministic state machine governing exception propagation, the precise conditions governing each transition, the bail-out trigger condition, the escalation algorithm, the bypass mechanism enforcing the no-machine-actor property, and the timeout handling that guarantees liveness under all failure conditions.

The protocol operates against the backdrop of the \bxthree{} three-layer model. The \purpose{} layer carries intent, judgment, and final authority over any decision exceeding a node's provisioned scope; it defines the operating range $R(n)$ for each node $n$ and receives all escalated exceptions. The \bounds{} layer carries bounded reasoning and constrained execution within $R(n)$; it evaluates whether proposed actions fall within the node's authority and initiates bailout when they do not. The \fact{} layer carries deterministic enforcement and forensic auditability; it verifies proposed actions against the authorization ledger, enforces the transition relation of the state machine, and maintains the append-only Forensic Ledger recording every protocol event as an immutable architectural record. Each spawned node carries an immutable parent pointer $\alpha(n)$, and every node carries a reference to its human accountability anchor $h(n)$ — a named human principal or institutional actor — established at node provisioning and never modified for the node's lifetime. The Bailout Protocol exploits this structure to guarantee that exception states propagate to $h(n)$ without passing through any intermediate machine actor as a terminus.

% ---------------------------------------------------------------
\subsection{State Machine Formalism}\label{sec:state-machine}

The Bailout Protocol is formally modeled as a deterministic finite state machine
\[
\mathcal{B} = (Q,\ q_0,\ \Sigma,\ \rightarrow,\ Q_F)
\]
where:

\begin{itemize}\itemsep0em
\item $Q = \{\,\text{IDLE},\ \text{BOUNDED},\ \text{BAIL\_PENDING},\ \text{ESCALATING},\ \text{HUMAN\_ACKNOWLEDGED},\ \text{RESOLVED}\,\}$ is the finite set of protocol states;
\item $q_0 = \text{IDLE}$ is the designated initial state, entered at node startup and after every successful directive resolution;
\item $\Sigma$ is the input alphabet, comprising the set of trigger events $e_{\text{trigger}}$, authorization confirmations $e_{\text{auth}}$, timeout signals $e_{\text{timeout}}$, human acknowledgments $e_{\text{ack}}$, binding resolution directives $e_{\text{resolve}}$, and rejection directives $e_{\text{reject}}$;
\item $\rightarrow\ \subseteq Q \times \Sigma \times Q$ is the deterministic transition relation (defined exhaustively in Table~\ref{tab:transition-relation});
\item $Q_F = \{\text{RESOLVED}\}$ is the set of terminal accepting states.
\end{itemize}

A node's state $q \in Q$ is stored in the node's worksheet and is observable by its parent node and by the Bailout Monitor. State updates are atomic; no intermediate or partially-executing states are observable by any external actor. The machine is \emph{deterministic}: for every $(q, \sigma) \in Q \times \Sigma$, there exists at most one $q' \in Q$ such that $(q, \sigma, q') \in \rightarrow$. Determinism is enforced by the Bailout Monitor's priority function $\pi$ (defined in \S\ref{sec:bailout-trigger}), which totally orders simultaneous trigger conditions and selects the highest-priority enabled transition before committing any state change. If no transition is defined for a $(q, \sigma)$ pair, the machine stalls — this is a protocol violation and is recorded as such in the Forensic Ledger.

\begin{table}[htbp]
\caption{Bailout Protocol state transition relation $\rightarrow\ \subseteq Q \times \Sigma \times Q$. Rows are current state; columns are input events. Entries denote the next state. Blank cells indicate that no transition is defined for that $(q, \sigma)$ pair — the machine stalls and the Bailout Monitor records a protocol violation.}
\label{tab:transition-relation}
\begin{tabular}{@{}lcccccc@{}}
\toprule
 & \multicolumn{6}{c}{Input event} \\
\cmidrule(l){2-7}
Current state & $e_{\text{trigger}}$ & $e_{\text{auth}}$ & $e_{\text{timeout}}$ & $e_{\text{ack}}$ & $e_{\text{resolve}}$ & $e_{\text{reject}}$ \\
\midrule
IDLE & BOUNDED & — & — & — & — & — \\
BOUNDED & — & HUMAN\_ACKNOWLEDGED & BAIL\_PENDING & — & — & — \\
BAIL\_PENDING & — & — & ESCALATING & — & — & — \\
ESCALATING & — & HUMAN\_ACKNOWLEDGED & ESCALATING & IDLE & RESOLVED & RESOLVED \\
HUMAN\_ACKNOWLEDGED & — & — & — & IDLE & RESOLVED & IDLE \\
RESOLVED & — & — & — & — & — & — \\
\bottomrule
\end{tabular}
\end{table}

\subparagraph{State definitions.}

\begin{itemize}\itemsep0em
\item[\textbf{IDLE.}] The node is operating within its provisioned authority. No exception condition is active. The Fact Layer processes normal action requests through the standard dispatch path. The node may transition to BOUNDED at any time upon receipt of a trigger event satisfying TC (defined in \S\ref{sec:bailout-trigger}).

\item[\textbf{BOUNDED.}] The node has detected an exception condition and is evaluating whether it is autonomously resolvable within current authority. The Bounds Engine is actively executing its resolution procedure. All machine-actor resolution paths remain open. If the Bounds Engine produces a resolution $a$ that the Fact Layer approves before $T_{\text{bounds}}$ expires, the node returns to IDLE via T2. If all resolution paths are exhausted or $T_{\text{bounds}}$ expires, the state advances to BAIL\_PENDING via T3.

\item[\textbf{BAIL\_PENDING.}] The node has determined — either through explicit \texttt{bailout\_signal} from the Bounds Engine or through $T_{\text{bounds}}$ expiration without Fact-Layer-approved resolution — that the exception exceeds its current authority. The Bailout Protocol is formally activated, but propagation to $h(n)$ has not yet been initiated. This is the last architectural point at which a machine actor could theoretically suppress escalation. It cannot: entry to BAIL\_PENDING triggers an immediate atomic transition to ESCALATING via T4, with no dwell time and no intermediate machine-actor intervention permitted.

\item[\textbf{ESCALATING.}] The exception is propagating upward through the parent chain toward $h(n)$. Each intermediate node receives the exception, appends its diagnostic context to the propagation chain (detailed in \S\ref{sec:escalation-algorithm}), and forwards the exception to its parent via \texttt{SendToParent} without attempting local resolution. The state remains ESCALATING until either (a) the human anchor acknowledges receipt ($e_{\text{ack}}$: T5), (b) the human anchor issues a binding directive ($e_{\text{resolve}}$ or $e_{\text{reject}}$: T7), or (c) all pathways time out after $N_{\max}$ retries at $T_{\text{cap}}$, propagating a \texttt{parent\_unreachable} event to $h(n)$ via an alternate functional pathway.

\item[\textbf{HUMAN\_ACKNOWLEDGED.}] The human anchor has received the exception, observed the diagnostic context, and acknowledged it. The node awaits a directive. No machine actor may issue, simulate, or substitute a directive in this state. The node remains in this state until the human principal issues \texttt{ACK} (T8: return to IDLE) or a binding directive \texttt{RESOLVE} / \texttt{REJECT} (T9: enter RESOLVED).

\item[\textbf{RESOLVED.}] The human principal has issued a directive and the protocol has recorded it in the authorization ledger. The directive is one of: \texttt{ACK} (continue with current parameters), \texttt{RESOLVE} (execute a specific binding resolution), or \texttt{REJECT} (disallow the proposed action and enter a quiescent error state). This state is terminal for the current protocol run. The node either returns to IDLE (on \texttt{ACK}) or becomes quiescent (on \texttt{REJECT}), awaiting human re-engagement.
\end{itemize}

\subparagraph{State invariants.} The following invariants hold for all reachable states of $\mathcal{B}$:

\begin{align}
\text{INV}_1 &: \forall n \in \text{nodes},\ \text{state}(n) \in Q \setminus \{\text{ESCALATING}\} \implies \neg\text{active\_bailout}(n). \tag{INV-1}
\end{align}
That is, a node occupies a non-escalating state \emph{iff} it has no active bailout propagation in flight. A node may carry a bailout history (a ledger entry recording a past propagation) while in IDLE, but no active propagation may be underway.

\begin{align}
\text{INV}_2 &: \text{state}(n) = \text{HUMAN\_ACKNOWLEDGED} \implies \text{directive\_issuer}(n) = h(n). \tag{INV-2}
\end{align}
Only the designated human anchor may issue directives in HUMAN\_ACKNOWLEDGED state. This invariant follows directly from the bypass mechanism (\S\ref{sec:bypass-mechanism}).

% ---------------------------------------------------------------
\subsection{Transition Conditions}\label{sec:transition-conditions}

Each transition in Table~\ref{tab:transition-relation} is governed by a precondition drawn from the node's operational context. We define each transition formally as a guarded command.

\begin{enumerate}
\item[\textbf{T1: IDLE $\rightarrow$ BOUNDED (trigger).}] Transition T1 fires when the node detects a \emph{bound signal} — any condition $x$ in the node's observable input space that falls outside the provisioned operating range $R(n)$ and satisfies the trigger condition TC. Formally:
\begin{equation}
\text{T1 fires} \iff \exists x \in \text{inputs}(n) : \text{TRIGGER}(n, x) \land \neg\text{resolved}(n, x). \tag{T1}
\end{equation}
The condition $\neg\text{resolved}(n, x)$ ensures that a previously resolved exception on the same input does not re-trigger. Resolution state is stored in the node's worksheet and is cleared only on system reset or explicit human override.

\item[\textbf{T2: BOUNDED $\rightarrow$ HUMAN\_ACKNOWLEDGED (intra-bounds resolution).}] Transition T2 fires when the Bounds Engine produces a resolution $a$ for exception $x$ that is approved by the Fact Layer before $T_{\text{bounds}}$ expires. Formally:
\begin{equation}
\text{T2 fires} \iff \exists a \in \text{resolutions}(n, x) : \text{FactLayer.approve}(a) \land \text{clock} < T_{\text{bounds}}. \tag{T2}
\end{equation}
If $T_{\text{bounds}}$ expires before a Fact-Layer-approved resolution is produced, the state transitions to BAIL\_PENDING via T3, not T2. The Fact Layer approval requirement ensures that even intra-bounds resolution — resolution achieved without human escalation — is audited and recorded.

\item[\textbf{T3: BOUNDED $\rightarrow$ BAIL\_PENDING (escalation).}] Transition T3 fires when $T_{\text{bounds}}$ expires without a Fact-Layer-approved resolution, or when the Bounds Engine emits an explicit \texttt{bailout\_signal} declaring that exception $x$ exceeds its current authority. Formally:
\begin{equation}
\text{T3 fires} \iff \bigl(\text{clock} \geq T_{\text{bounds}} \land \neg\exists a : \text{FactLayer.approve}(a)\bigr) \;\lor\; \texttt{bailout\_signal}(x). \tag{T3}
\end{equation}
The disjunction reflects two independent pathways into BAIL\_PENDING: timeout-driven (the Bounds Engine failed to resolve within the allotted time) and signal-driven (the Bounds Engine explicitly declared the exception beyond its scope).

\item[\textbf{T4: BAIL\_PENDING $\rightarrow$ ESCALATING (protocol activation).}] Transition T4 fires immediately upon entry to BAIL\_PENDING — there is no dwell time, no retry attempt, and no machine-actor suppression option. The Bailout Monitor atomically commits the BAIL\_PENDING entry to the Forensic Ledger and immediately invokes the \texttt{Escalate} algorithm (\S\ref{sec:escalation-algorithm}) to initiate propagation to $h(n)$. Formally:
\begin{equation}
\text{T4 fires} \iff \text{state} = \text{BAIL\_PENDING}. \tag{T4}
\end{equation}

\item[\textbf{T5: ESCALATING $\rightarrow$ HUMAN\_ACKNOWLEDGED (human directive received).}] Transition T5 fires when the human anchor $h(n)$ issues a directive $d \in \{\texttt{ACK},\ \texttt{RESOLVE},\ \texttt{REJECT}\}$ and the Fact Layer pre-commit hook verifies the directive originated from $h(n)$ (not from any machine actor). Formally:
\begin{equation}
\text{T5 fires} \iff d \in \{\texttt{ACK},\ \texttt{RESOLVE},\ \texttt{REJECT}\} \land \text{FactLayer.verifyOrigin}(d, h(n)). \tag{T5}
\end{equation}
If the directive is \texttt{ACK}, the node returns to IDLE (T8). If it is \texttt{RESOLVE} or \texttt{REJECT}, the protocol enters RESOLVED (T9).

\item[\textbf{T6: ESCALATING $\rightarrow$ ESCALATING (retry on timeout).}] Transition T6 fires when the propagation timeout $T_{\text{prop}}$ expires without receiving an acknowledgment from the parent node. The algorithm retries with exponential backoff. Formally:
\begin{equation}
\text{T6 fires} \iff \text{clock} \geq T_{\text{prop}} \land \text{attempts} < N_{\max}. \tag{T6}
\end{equation}

\item[\textbf{T7: ESCALATING $\rightarrow$ RESOLVED (human timeout).}] Transition T7 fires when all $N_{\max}$ retry attempts have been exhausted without receiving acknowledgment from the parent. A \texttt{parent\_unreachable} event is recorded in the Forensic Ledger and propagated to $h(n)$ via the next available functional pathway. Formally:
\begin{equation}
\text{T7 fires} \iff \text{attempts} \geq N_{\max} \land \neg\text{acknowledged}. \tag{T7}
\end{equation}
\end{enumerate}

% ---------------------------------------------------------------
\subsection{Bail-Out Trigger Condition}\label{sec:bailout-trigger}

The bail-out trigger is the condition that causes a node to enter BOUNDED state (T1). Rather than requiring system designers to enumerate specific error codes in advance — a practice that provides coverage only for anticipated failure modes and leaves unanticipated ones invisible — the Bailout Protocol defines the trigger in terms of \emph{observational boundary violation}: a node escalates precisely when its observable input space contains a condition outside its provisioned operating range $R(n)$, and the node cannot produce a Fact-Layer-approved resolution within $T_{\text{bounds}}$.

\subparagraph{Primary trigger condition.} The trigger $\text{TRIGGER}(n, x)$ is defined as a conjunction of three sub-conditions, all of which must hold simultaneously:
\begin{equation}
\text{TRIGGER}(n, x) \;\equiv\; \underbrace{x \notin R(n)}_{(i)} \;\land\; \underbrace{\neg\text{resolved}(n, x)}_{(ii)} \;\land\; \underbrace{\text{confidence}(n, x) < \tau}_{(iii)} \tag{TC}
\end{equation}
where:

\begin{itemize}\itemsep0em
\item[(i)] $x$ is an element of the node's observable input space that falls outside the provisioned operating range $R(n)$. This is the \emph{bound signal} condition: the node observes something it was not provisioned to handle autonomously.
\item[(ii)] $x$ has not been previously resolved by the node within the current operational session. This prevents re-trigger on known-exception inputs that have already received human adjudication.
\item[(iii)] $\text{confidence}(n, x) < \tau$ is the node's self-assessed confidence that it can produce a correct resolution for $x$ without human input, where $\tau$ is the provisioned confidence threshold. The confidence function $\text{confidence} : \text{node} \times \text{input} \rightarrow [0, 1]$ is defined at node provisioning and is \emph{not} adjustable at runtime by any machine actor. This prevents a machine actor from inflating its own confidence to suppress a warranted escalation.
\end{itemize}

\subparagraph{Trigger priority function.} When multiple trigger conditions fire simultaneously for node $n$, they are totally ordered by the priority function $\pi(n, x)$, which ensures deterministic transition selection:
\begin{equation}
\pi(n, x) \;=\; w_1 \cdot \text{severity}(x) \;+\; w_2 \cdot \bigl(1 - \text{confidence}(n, x)\bigr) \;+\; w_3 \cdot \text{distance}\bigl(n, h(n)\bigr) \tag{PF}
\end{equation}
where $\text{severity}(x) \in [0, 1]$ is the assigned severity of exception $x$ (higher values indicate greater potential impact), $\text{distance}(n, h(n))$ is the number of hops from $n$ to the human anchor in the parent chain (longer distances increase priority because propagation latency is greater), and $w_1, w_2, w_3$ are provisioned constants satisfying $w_1 > w_2 > w_3 \geq 0$. The highest-priority trigger fires first. Ordering is enforced by the Bailout Monitor before any state transition is committed to the Forensic Ledger.

\subparagraph{Secondary trigger: explicit \texttt{bailout\_signal}.} In addition to TC, the Bounds Engine may emit an explicit \texttt{bailout\_signal} at any time, regardless of confidence scores. This mechanism handles the case where the Bounds Engine detects a condition — for example, a logical contradiction between two inputs within $R(n)$ — that it can recognize as unresolvable even though both inputs are individually within the operating range. The \texttt{bailout\_signal} bypasses TC entirely and immediately initiates T3:
\begin{equation}
\texttt{bailout\_signal}(x) \;\implies\; \text{T3 fires for node } n \text{ with exception } x. \tag{TC-secondary}
\end{equation}

% ---------------------------------------------------------------
\subsection{Escalation Algorithm}\label{sec:escalation-algorithm}

When the state machine enters ESCALATING, the protocol executes the \texttt{Escalate} algorithm to propagate the exception to $h(n)$ along the parent chain. The algorithm is executed by the Bailout Monitor on the originating node; intermediate nodes execute only \texttt{ForwardException}, which appends diagnostic context and relays the exception upward.

\begin{algorithm}[htbp]
\caption{Bailout Protocol escalation algorithm.}
\label{alg:escalate}
\begin{algorithmic}[1]
\Procedure{Escalate}{$n$, trigger}
  \State $\text{ctx} \leftarrow$ \Call{BuildDiagnosticContext}{$n$, trigger}
  \State $p \leftarrow \alpha(n)$
  \State $\text{attempts} \leftarrow 0$
  \State $T_{\text{prop}} \leftarrow T_{\text{initial}}$
  \Repeat
    \State $\text{pathway} \leftarrow$ \Call{SelectPathway}{PHR, $p$}
    \State $\text{acked} \leftarrow$ \Call{SendToParent}{$p$, ctx, $T_{\text{prop}}$, pathway}
    \State \Call{UpdateLedger}{\text{propagation\_attempt}, $n$, $p$, $\text{pathway}$, $\text{clock}$}
    \If{$\text{acked}$}
      \State \Call{Transition}{\text{HUMAN\_ACKNOWLEDGED}}
      \State \Return \texttt{propagated}
    \EndIf
    \State $\text{attempts} \leftarrow \text{attempts} + 1$
    \State $T_{\text{prop}} \leftarrow \min(2 \cdot T_{\text{prop}},\ T_{\text{cap}})$
    \State \Call{UpdateLedger}{\text{retry}, \text{attempts}, \text{clock}$}
  \Until $\text{attempts} \geq N_{\max}$
  \State \Call{UpdateLedger}{\text{parent\_unreachable}, n, p, \text{clock}$}
  \State \Call{Transition}{\text{RESOLVED}}
  \State \Return \texttt{unresolved}
\EndProcedure
\\[6pt]
\Procedure{BuildDiagnosticContext}{$n$, trigger}
  \State $\text{ctx} \leftarrow \emptyset$
  \State $\text{ctx.trigger} \leftarrow trigger$
  \State $\text{ctx.node\_id} \leftarrow n$
  \State $\text{ctx.parent\_chain} \leftarrow ()$
  \State $\text{ctx.timestamp} \leftarrow \text{clock}()$
  \State $\text{ctx.propagation\_chain} \leftarrow []$
  \State \Return ctx
\EndProcedure
\\[6pt]
\Procedure{SelectPathway}{PHR, target}
  \State $\text{candidates} \leftarrow$ PHR.$\text{functional\_pathways\_to}$(target)
  \State \Return $\displaystyle\arg\min_{\text{pht}\in\text{candidates}} \text{pht.estimated\_latency}$
\EndProcedure
\end{algorithmic}
\end{algorithm}

\subparagraph{Propagation chain integrity.} The diagnostic context $\text{ctx}$ is immutable after construction by \texttt{BuildDiagnosticContext}. No node along the propagation path may modify, truncate, redact, or suppress any field of $\text{ctx}$. Each node appends a record to $\text{ctx.propagation\_chain}$ containing: its node identifier, a timestamp, a state snapshot, and any locally relevant diagnostic metadata. This append-only integrity property is enforced by the Fact Layer: the \texttt{UpdateLedger} call is atomic, and any attempt to modify an existing ledger entry is rejected and logged as a protocol violation.

\subparagraph{Correctness properties of \texttt{Escalate}.}

\begin{enumerate}
\item[\emph{Progress.}] If the parent node $p$ is reachable via at least one functional pathway, the algorithm eventually receives an ACK and returns \texttt{propagated}. The retry loop implements exponential backoff capped at $T_{\text{cap}}$, guaranteeing termination within a bounded number of attempts $N_{\max}$. If $p$ is unreachable via all provisioned pathways for $N_{\max}$ consecutive attempts, the algorithm fires a \texttt{parent\_unreachable} timeout event (T7), which propagates to $h(n)$ via the next available functional pathway tracked by the Pathway Health Registry.

\item[\emph{Integrity.}] The diagnostic context is append-only and immutable after construction. No machine actor can alter, suppress, or selectively omit fields of the propagation chain. The Fact Layer enforces this property at the ledger level: every \texttt{UpdateLedger} call is recorded, and the chain of records is cryptographically sealed (via SHA-256, as specified in the \bxthree{} Forensic Ledger specification).

\item[\emph{Liveness under pathway degradation.}] If $p$ is unreachable, the algorithm retries with exponential backoff. After $N_{\max}$ failures at $T_{\text{cap}}$, it fires a \texttt{parent\_unreachable} event, which propagates to $h(n)$ via the next available functional pathway in the PHR. Pathways that fail $K$ consecutive health checks are marked \texttt{DEGRADED} and excluded from pathway selection until a health check succeeds, preventing the algorithm from repeatedly selecting a known-failed pathway.
\end{enumerate}

% ---------------------------------------------------------------
\subsection{Bypass Mechanism}\label{sec:bypass-mechanism}

The bypass mechanism is the structural property that distinguishes the Bailout Protocol from conventional escalation hierarchies. It ensures that no machine actor — regardless of its position in the parent chain, its capability level, or its access to the Fact Layer — can absorb, resolve, or suppress an exception that has entered the ESCALATING state.

\subparagraph{Formal definition.} The bypass property is stated formally as:
\begin{equation}
\text{BYPASS}(n) \;\equiv\; \forall m \in \text{MachineActors} \setminus \{h(n)\}:\; m \notin \text{terminus}\bigl(\text{ESCALATING}(n)\bigr) \tag{BP}
\end{equation}
where $\text{terminus}(\text{ESCALATING}(n))$ denotes the set of actors that may be the final recipient of the ESCALATING state for node $n$. BYPASS requires that this set contains exactly one element: the human anchor $h(n)$. No machine actor may appear in this set.

\subparagraph{Why bypass is architecturally necessary.} The bypass is not an optimization — it is a structural necessity that follows from the upstream accountability guarantee of the \bxthree{} Framework. Any machine actor $m$ that could absorb an exception in ESCALATING state would become a terminal node for that exception, and the accountability chain would terminate at $m$ rather than at a human. If $m$ then fails, or if $m$'s resolution is incorrect or insufficient, no human principal is present in the chain to detect or correct the failure. The system compounds failure instead of surfacing it. This is the precise failure mode that the Bailout Protocol is designed to prevent.

Conventional escalation hierarchies — including the Tiered Agentic Oversight (TAO) model — permit this compounding by making human oversight a high-tier option at the top of a machine-tier hierarchy, where each intermediate tier may absorb the exception before passing it up. The human anchor in TAO is present but not architecturally required for every exception; absorption at an intermediate tier is permitted and invisible to external auditors. The Bailout Protocol forbids this absorption by architectural constraint: the human anchor is not one option among many — it is the \emph{only} permissible terminus for every unresolved exception. The parent-chain bypass skips intermediate nodes not solely for efficiency but because their presence in the terminus would break the accountability chain by introducing machine actors as potential terminal nodes.

\subparagraph{Enforcement mechanisms.} The bypass is enforced by two cooperating mechanisms that operate at different layers of the \bxthree{} architecture:

\begin{enumerate}
\item[\textbf{Fact Layer gate (enforcement at the authorization layer).}] The Fact Layer is the only mechanism capable of authorizing state transitions in the \bxthree{} Framework. When a node $n$ is in BOUNDED state and produces a resolution $a$, the resolution is written to the authorization ledger only after Fact Layer approval. The Bailout Monitor's pre-commit hook rejects any ledger entry for a resolution where the underlying trigger satisfies TC and $T_{\text{bounds}}$ has expired. This prevents the Bounds Engine from issuing a retroactive resolution after timeout to suppress warranted escalation. The Fact Layer gate operates identically on all machine actors: there is no privileged machine actor with the ability to bypass this check.

\item[\textbf{Immutable parent-pointer and human-root architecture (enforcement at the data model layer).}] The parent pointer $\alpha(n)$ and the human-root identifier $h(n)$ are stored in Fact Layer-managed memory and are read-only at the machine-actor level. A machine actor cannot redirect the parent chain, cannot insert itself as an intermediate resolver, and cannot modify the $humanRootId$ field that terminates the chain. Any attempted modification requires Purpose Layer authorization with a corresponding ledger entry — which cannot be completed without human approval if the modification would affect the bailout terminus. This architectural immutability ensures that the propagation path is fixed at provisioning and cannot be altered by any machine actor at runtime.
\end{enumerate}

Together, these two mechanisms ensure that BYPASS holds for all $n \in \text{nodes}$ under all execution conditions, including partial node failure, network partition, and deliberate attempt by a compromised machine actor to redirect the bailout pathway.

% ---------------------------------------------------------------
\subsection{Timeout Handling}\label{sec:timeout-handling}

Timeouts in the Bailout Protocol serve a liveness function: they prevent the protocol from permanently stalling when a parent node is unreachable, a communication pathway is unresponsive, or the human anchor is temporarily unavailable. Each timeout has a defined scope, a defined scope-exit action, and an escalation path that prevents the protocol from deadlocking.

\subparagraph{Timeout definitions.}

\begin{enumerate}
\item[\textbf{$T_{\text{bounds}}$ — Bounds Engine resolution timeout.}] The maximum time a node may remain in BOUNDED state attempting autonomous resolution. Default: 60 seconds, provisioned per node class according to the expected latency of the node's resolution procedure. If $T_{\text{bounds}}$ expires without a Fact-Layer-approved resolution, the node transitions to BAIL\_PENDING $\rightarrow$ ESCALATING via T3 and T4. The Fact Layer cannot extend $T_{\text{bounds}}$ unilaterally; any extension requires human authorization via a separate ledger entry recorded before $T_{\text{bounds}}$ expires.

\item[\textbf{$T_{\text{prop}}$ — Per-hop propagation timeout.}] The maximum time node $n$ waits for an acknowledgment from its parent $p$ after sending an exception via \texttt{SendToParent}. Default initial value: 5 seconds. If $T_{\text{prop}}$ expires without ACK, the algorithm retries with exponential backoff: $T_{\text{prop}} \leftarrow \min(2 \cdot T_{\text{prop}},\ T_{\text{cap}})$. This backoff continues until either an ACK is received or $N_{\max}$ retries have been exhausted at $T_{\text{cap}}$.

\item[\textbf{$T_{\text{cap}}$ — Maximum propagation timeout value.}] The ceiling value for the exponential backoff of $T_{\text{prop}}$. Default: 60 seconds. After $N_{\max}$ consecutive retries at $T_{\text{cap}}$ without ACK, the algorithm fires a \texttt{parent\_unreachable} timeout event (T7), which is treated as a new trigger event propagating to $h(n)$ via the next available functional pathway tracked by the Pathway Health Registry.

\item[\textbf{$T_{\text{human}}$ — Human acknowledgment timeout.}] The maximum time the human anchor may remain in HUMAN\_ACKNOWLEDGED state without issuing a directive. Default: 300 seconds (5 minutes), configurable per deployment based on the expected human response time for the node's operational domain. If $T_{\text{human}}$ expires, the protocol issues a reminder notification to $h(n)$ and transitions the originating node to RESOLVED with an \texttt{unresolved} tag. The node becomes quiescent: no autonomous retry occurs, no further actions are attempted, and the \texttt{unresolved} tag in the Forensic Ledger signals that human re-engagement is required to clear the state. The protocol does not attempt to re-escalate automatically; automatic re-escalation would create a potential denial-of-service condition against the human anchor.
\end{enumerate}

\begin{figure}[htbp]
\centering
\begin{tabular}{lcl}
\hline
\textbf{Timeout} & \textbf{Started at} & \textbf{Scope} \\ \hline
$T_{\text{bounds}}$ & T1 fires (TC) & Bounds Engine autonomous resolution window \\
$T_{\text{prop}}$ & T4 fires (transmit to $h(n)$) & Per-hop parent acknowledgment \\
$T_{\text{cap}}$ & T6 fires (retry) & Ceiling for $T_{\text{prop}}$ backoff \\
$T_{\text{human}}$ & ESCALATING + contact confirmed & Human directive issuance window \\
\hline
\end{tabular}
\caption{Timeout parameters, their start conditions, and their functional scopes.}
\label{tab:timeouts}
\end{figure}

\subparagraph{Pathway Health Registry (PHR).} The PHR is a runtime data structure maintained by the Bailout Monitor that tracks the availability, latency, and error rate of each provisioned communication pathway to the human anchor. Before each \texttt{SendToParent} call, \texttt{SelectPathway} queries the PHR and selects the pathway with the lowest estimated latency among currently functional pathways. Pathways that fail $K$ consecutive health checks (default $K = 3$) are marked \texttt{DEGRADED} and excluded from pathway selection until a health check succeeds, preventing the algorithm from repeatedly selecting a known-degraded pathway.

The PHR is updated by a background health-ping thread that sends heartbeat messages along each provisioned pathway at intervals of $T_{\text{health}}$ (default: 30 seconds). The health-ping thread operates with the same scheduling priority as the Bailout Monitor and is not preemptible by agent tasks, ensuring that pathway status remains current even under heavy agent workload.

\subparagraph{Liveness guarantee.} Together, $T_{\text{bounds}}$, $T_{\text{prop}}$, $T_{\text{cap}}$, and $T_{\text{human}}$ define a finite upper bound on the wall-clock time from trigger firing (T1) to human acknowledgment under any combination of parent unavailability and pathway degradation. This bound is a deployment parameter — not a protocol constant — derived from the network topology, the number of hops to the human anchor, and the provisioned $T_{\max}$ parameter. A deployment that cannot guarantee human acknowledgment within the required $T_{\max}$ under all failure scenarios must provision additional or higher-bandwidth pathways, or must reduce $T_{\text{prop}}$ and $T_{\text{cap}}$ to bring the worst-case bound within the required limit.
